
donc tel que @@M0@@ soit un @@M1@@-Module de type fini, donc de présentation finie d'après les hypothèses noethériennes; le résultat est d'ailleurs valable pour tout entier @@M2@@, les @@M3@@ étant cohérents d'après le théorème de Serre.

La proposition est donc montrée dans le cas où @@M4@@ est noethérien.

Supposons maintenant qu'on ait un carré cartésien :
@@M5@@
avec @@M6@@ noethérien, @@M7@@; soit @@M8@@ un @@M9@@-Module de présentation finie, plat sur @@M10@@, et soit @@M11@@. D'après ce qui précède la proposition est vraie pour @@M12@@. D'après EGA III 7.7.5 et 7.8.5, le théorème de changement de base est vrai avec les hypothèses précédentes; on en déduit aussitôt que le théorème est vrai pour @@M13@@.

Il suffit donc, pour achever la démonstration, de montrer que le cas général ( avec @@M14@@ affine) peut être obtenu à partir du cas noethérien par changement de base.

Comme @@M15@@ est affine, @@M16@@ est limite projective filtrante de schémas affines noethériens (si @@M17@@, on écrit @@M18@@ comme limite inductive de ses sous-algèbres de type fini sur @@M19@@). Soit donc @@M20@@. D'après EGA IV 8.5.2 et 8.5.5, il existe un indice @@M21@@ et un @@M22@@-Module localement libre de type fini @@M23@@ tel que @@M24@@ (en notant @@M25@@ le morphisme @@M26@@, et @@M27@@ les morphismes de transition @@M28@@).

Si on pose, pour @@M29@@, @@M30@@, et @@M31@@, on a pour tout @@M32@@, @@M33@@, et @@M34@@. Par suite, il existe un indice @@M35@@
